\documentclass[11pt,a4paper]{article}
\usepackage[utf8]{inputenc}
\usepackage[margin=1in]{geometry}
\usepackage{amsmath,amssymb}
\usepackage{booktabs}
\usepackage{hyperref}
\usepackage{xcolor}
\usepackage{enumitem}

\definecolor{navy}{RGB}{0,32,96}
\definecolor{darkblue}{RGB}{0,51,153}
\definecolor{graybox}{RGB}{245,245,248}

\hypersetup{
    colorlinks=true,
    linkcolor=darkblue,
    citecolor=darkblue,
    urlcolor=darkblue
}

\begin{document}

\begin{center}
    {\Large \textbf{\textcolor{navy}{Point-by-Point Author Response to Reviewers}}}\\[6pt]
    \textbf{Manuscript ID: IDCIoT-154}\\[3pt]
    \textbf{Title:} CoT-HMM: Structural Plausibility Scoring of Chain-of-Thought Reasoning via Hidden Markov Models\\[3pt]
    \textbf{Authors:} Gaurav Patil, Chirayu Marathe, Parth Parmar, Suhail Shaikh, Bhagyalakshmi V\\[3pt]
    \textit{5th International Conference on Intelligent Data Communication Technologies and Internet of Things (IDCIoT 2027)}\\[2pt]
    \textit{IEEE Conference Record \#71706}
\end{center}

\vspace{10pt}
\hrule
\vspace{12pt}

We sincerely thank the Technical Program Committee and both anonymous reviewers for their constructive, rigorous, and encouraging feedback on our manuscript. We have thoroughly revised the paper according to your recommendations. Below is our detailed point-by-point response addressing each remark, with direct references to the revised sections of the manuscript.

\section*{Response to Reviewer 1}

\noindent\textbf{Remark R1.1:} \textit{Evaluates only 39 correct traces from Qwen2.5-0.5B; sensitive sample size.}
\begin{itemize}[leftmargin=18pt]
    \item \textbf{Response:} In Section IV-A (\textit{Setup: Baselines and Metrics}), we have explicitly clarified that the $N=39$ verified-correct traces from Qwen2.5-0.5B served strictly as a focused, high-precision seed pilot corpus for controlled structural error injection.
    \item To guard against small-sample variance, non-parametric percentile bootstrap resampling ($B=10{,}000$ iterations) was used to compute robust 95\% confidence intervals across all reported metrics (Table II).
    \item Furthermore, in Section IV-C and Table IV (\textit{Zero-Shot Cross-Architecture Transfer}), we demonstrate that the transition grammar learned from this seed corpus transfers zero-shot to substantially larger model families without retraining: Llama-3.2-1B ($r=0.946$), Phi-3-mini ($r=0.941$), and Mistral-7B ($r=0.970$) (all $p < 0.001$), confirming generalizability across model scales.
\end{itemize}

\vspace{6pt}
\noindent\textbf{Remark R1.2:} \textit{Synthetic corruption vs. natural reasoning failures.}
\begin{itemize}[leftmargin=18pt]
    \item \textbf{Response:} Controlled perturbation (step deletion, loop injection, premature truncation) enables causal isolation of structural failure modes without confounding noise from arithmetic mistakes or lexical vocabulary shifts, an established methodology in reasoning evaluation benchmarks (e.g., ROSCOE).
    \item Crucially, this synthetic detection directly mirrors real-world reasoning failures: in uncorrupted, organic model generations on GSM8K (Section IV-D), \textbf{77.8\% of looping chains ($N=150$) were mathematically incorrect vs.\ 23.8\% of non-looping ones} ($z=6.18, p=6.3\times 10^{-10}$).
\end{itemize}

\vspace{6pt}
\noindent\textbf{Remark R1.3:} \textit{Regex-based semantic tagger ambiguity \& error propagation.}
\begin{itemize}[leftmargin=18pt]
    \item \textbf{Response:} In Section III-B, we report that our rule-based tagger achieved Cohen's $\kappa=0.953$ agreement against a 50-step expert gold standard (comparable to Phi-3-mini's $\kappa=0.976$).
    \item We have added a dedicated paragraph (\textit{Tagger Error Analysis}) in Section III-B detailing failure modes: multi-clause sentences with overlapping discourse markers (e.g., ``Hence $x=2$'') and colloquial discourse markers. Because the Ensemble HMM models continuous emission distributions across $M=12$ models, isolated tagging fluctuations are softly absorbed across the full sequence. In Section VI, fine-tuning a lightweight DeBERTa-v3 sentence classifier is proposed as a direct future extension.
\end{itemize}

\vspace{6pt}
\noindent\textbf{Remark R1.4:} \textit{Table II AUC: 0.658 for missing computation \& 0.761 for circular reasoning vs. high for premature conclusions.}
\begin{itemize}[leftmargin=18pt]
    \item \textbf{Response:} In Section IV-B, we added an explanatory paragraph (\textit{Asymmetric Detection Mechanism}):
    \begin{itemize}
        \item \textit{Premature Conclusions (AUC = 0.990):} Directly jumping from problem restatement to conclusion ($\text{PR}\to\text{CO}$) represents a severe terminal transition violation ($P(\text{CO}\mid\text{PR}) < 10^{-4}$ in the learned matrix), triggering a catastrophic log-likelihood penalty ($d=4.166$).
        \item \textit{Missing Computation (AUC = 0.658):} Omitting arithmetic calculation tags ($\text{CA}$) leaves residual sequences of logical deductions ($\text{PR}\to\text{IN}\to\text{IN}\to\text{CO}$) that locally resemble valid qualitative reasoning, softening the likelihood drop. We note that token-length weighting and entropy thresholds provide immediate avenues to boost missing-step sensitivity.
    \end{itemize}
\end{itemize}

\vspace{6pt}
\noindent\textbf{Remark R1.5:} \textit{Distinction between Table II (0.990) and Table III (0.982) AUC.}
\begin{itemize}[leftmargin=18pt]
    \item \textbf{Response:} Table II reports \textit{per-category controlled error injection AUCs} on the pilot corpus to isolate specific failure modes. Table III reports the \textit{composite out-of-fold plausibility discrimination AUC} ($0.982$) on mixed benchmark traces against baseline verifiers (ROSCOE, PiCSAR, Length, Random). We added clear cross-referencing in Section IV-B to clarify this distinction.
\end{itemize}

\vspace{6pt}
\noindent\textbf{Remark R1.6:} \textit{Conclusion should be a separate section after Results.}
\begin{itemize}[leftmargin=18pt]
    \item \textbf{Response:} We restructured the manuscript to separate Discussion and Conclusion. Section V is now dedicated to \textbf{Discussion}, and Section VI is a standalone section titled \textbf{Conclusion and Future Work}.
\end{itemize}

\vspace{6pt}
\noindent\textbf{Remark R1.7:} \textit{Proofread for grammar/spelling/typos.}
\begin{itemize}[leftmargin=18pt]
    \item \textbf{Response:} A thorough editorial pass was performed across all sections, refining notation, equations, and caption clarity.
\end{itemize}

\newpage
\section*{Response to Reviewer 2}

\noindent\textbf{Remark R2.1:} \textit{How does CoT-HMM handle arithmetic errors?}
\begin{itemize}[leftmargin=18pt]
    \item \textbf{Response:} In Section V, we added a dedicated subsection titled \textbf{Structural vs.\ Arithmetic Correctness}. CoT-HMM evaluates reasoning grammar, not numerical arithmetic. An equation such as $12 \times 4 = 52$ possesses valid reasoning syntax ($\text{EQ}\to\text{CA}\to\text{IN}$); hence, CoT-HMM appropriately scores it as structurally plausible.
    \item This is empirically demonstrated in Table II, where our semantic-only corruption condition (pristine structure, incorrect arithmetic) yielded AUC = $0.501$ (random guessing). Detecting numerical errors belongs to downstream deterministic tools (e.g., Python REPL / calculators). CoT-HMM acts as a fast ($\approx$150ms), zero-GPU filter that eliminates structurally broken chains before invoking heavier tools.
\end{itemize}

\vspace{6pt}
\noindent\textbf{Remark R2.2:} \textit{Missing-computation detection AUC remains insufficient (0.658).}
\begin{itemize}[leftmargin=18pt]
    \item \textbf{Response:} We concur with the reviewer. As explained in Remark R1.4, omitting calculation steps preserves valid qualitative deduction sub-chains. In Section IV-B and Section VI, we explicitly discuss this trade-off and outline mitigations via token-length weighted transition matrices and posterior Structural Entropy thresholds.
\end{itemize}

\vspace{6pt}
\noindent\textbf{Remark R2.3:} \textit{How was the 99.3\% routing accuracy calculated?}
\begin{itemize}[leftmargin=18pt]
    \item \textbf{Response:} In Section IV-D, we have included the explicit mathematical formula:
    \begin{equation}
    \text{Acc}_{\text{pass}} = \frac{|\mathcal{T}_{\text{pass}} \cap \mathcal{C}_{\text{correct}}|}{|\mathcal{T}_{\text{pass}}|} = \frac{TP}{TP + FP} = 99.3\%
    \end{equation}
    where $\mathcal{T}_{\text{pass}}$ contains traces with plausibility margin exceeding the decision threshold $\tau_{\text{high}}$ (calibrated via Youden's $J$ on validation data), and $\mathcal{C}_{\text{correct}}$ denotes ground-truth correct chains. Traces in this top 35\% tier bypass secondary verification with 99.3\% accuracy.
\end{itemize}

\vspace{6pt}
\noindent\textbf{Remark R2.4:} \textit{Regex-based semantic tagging needs error analysis.}
\begin{itemize}[leftmargin=18pt]
    \item \textbf{Response:} Addressed in Section III-B (\textit{Tagger Error Analysis}), outlining failure modes (overlapping discourse markers, conversational cues) and explaining why the Ensemble HMM's soft emission distributions prevent individual misclassifications from corrupting sequence-level plausibility scores.
\end{itemize}

\vspace{6pt}
\noindent\textbf{Remark R2.5:} \textit{Semantic error detection limitations require discussion.}
\begin{itemize}[leftmargin=18pt]
    \item \textbf{Response:} In Section V (\textit{Discussion}), under \textbf{Structural vs.\ Arithmetic Correctness} and \textbf{Limitations}, we expand on the complementary division of labor: CoT-HMM rapidly detects macroscopic structural breakdowns (loops, premature stops, disordered logic), while specialized semantic verifiers (e.g., PRMs, Python REPL) verify numerical and factual correctness, reducing total token overhead by 42\%--47\%.
\end{itemize}

\vspace{12pt}
\begin{center}
\textbf{Summary of Manuscript Modifications (IEEE IDCIoT-154)}
\vspace{4pt}

\begin{tabular}{lp{10cm}}
\toprule
\textbf{Section} & \textbf{Revisions Incorporated} \\
\midrule
Section III-B & Added tagger error analysis, regex failure modes, and neural tagger roadmap. \\
Section IV-A & Clarified $N=39$ pilot corpus, $B=10{,}000$ bootstrap CIs, and multi-model transfer. \\
Section IV-B & Justified synthetic perturbation; detailed asymmetric AUC mechanism; cross-referenced Table II vs.\ Table III. \\
Section IV-D & Added Eq.~(3) for 99.3\% pass-tier compute routing accuracy. \\
Section V & Dedicated Discussion section with \textit{Structural vs.\ Arithmetic Correctness}. \\
Section VI & Created standalone \textit{Conclusion and Future Work} section. \\
Throughout & Thorough proofreading and verified exact 5-page IEEEtran layout compliance. \\
\bottomrule
\end{tabular}
\end{center}

\end{document}
